\section*{Hoofdstuk \ref{ch:g4:journey-of-food} --- De reis van het voedsel}

\begin{solution}{exo:g4:journey-of-food:1}
\omterm{def:g4:journey-of-food:tract}{Mond}, \omterm{def:g4:journey-of-food:tract}{slokdarm}, \omterm{def:g4:journey-of-food:tract}{maag}, \omterm{def:g4:journey-of-food:tract}{dunne darm}, \omterm{def:g4:journey-of-food:tract}{dikke darm} --- en daarna verlaten de
resten het lichaam.
\end{solution}

\begin{solution}{exo:g4:journey-of-food:2}
Het wordt door de gespierde wanden van de \omterm{def:g4:journey-of-food:tract}{maag} gekneed en met
\omterm{def:g4:journey-of-food:digestion}{verteringssappen} gemengd --- gedurende enkele uren.
\end{solution}

\begin{solution}{exo:g4:journey-of-food:3}
\omterm{def:g4:journey-of-food:digestion}{Vertering} is de omzetting van voedsel, langs het kanaal, in
bestanddelen die eenvoudig en klein genoeg zijn om het lichaam binnen
te gaan. Die bestanddelen zijn de \omterm{def:g4:journey-of-food:digestion}{voedingsstoffen}.
\end{solution}

\begin{solution}{exo:g4:journey-of-food:4}
In de \omterm{def:g4:journey-of-food:tract}{dunne darm}, door haar dunne geplooide wanden, het \omterm{prop:g4:journey-of-food:absorb}{bloed} in. Wat
niet kan oversteken reist verder naar de \omterm{def:g4:journey-of-food:tract}{dikke darm}, verliest daar zijn
water en verlaat het lichaam als ontlasting.
\end{solution}

\begin{solution}{exo:g4:journey-of-food:5}
Het \omterm{def:g4:journey-of-food:tract}{speeksel} is het brood al in eenvoudiger, zoetere bestanddelen aan
het afbreken --- de \omterm{def:g4:journey-of-food:digestion}{vertering} begint in de \omterm{def:g4:journey-of-food:tract}{mond}, en de zoetheid is haar
smaak.
\end{solution}

\begin{solution}{exo:g4:journey-of-food:6}
Omdat oversteken naar het \omterm{prop:g4:journey-of-food:absorb}{bloed} wandoppervlak kost: de meters lengte en
de diepe plooien vanbinnen persen een enorm oversteekoppervlak ---
uitgespreid ongeveer het tapijt van een kleine kamer --- in de buik.
\end{solution}

\begin{solution}{exo:g4:journey-of-food:7}
Bijvoorbeeld: goed kauwen (de \omterm{def:g4:journey-of-food:tract}{maag} heeft geen tanden); rustig en op
vaste tijden eten (het kanaal werkt het best ongehaast); water drinken
(de \omterm{def:g4:journey-of-food:digestion}{vertering} loopt nat); vezels eten (die houden de resten in
beweging); na grote maaltijden rusten (\omterm{def:g4:journey-of-food:tract}{maag} en \omterm{def:g1:my-body:parts}{benen} concurreren om de
inspanning van het lichaam).
\end{solution}

\begin{solution}{exo:g4:journey-of-food:8}
Open antwoord. Meestal ruim onder het halfuur in de \omterm{def:g4:journey-of-food:tract}{mond} --- tegenover
een tot twee dagen voor de rest van de reis. Het aandeel van de \omterm{def:g4:journey-of-food:tract}{mond} is
klein in tijd en beslissend in uitwerking: het is de enige halte met
tanden.
\end{solution}

\begin{solution}{exo:g4:journey-of-food:9}
Gebroken: goed kauwen (de lunch werd halfgekauwd naar binnen geschrokt)
en na de maaltijd rusten (meteen daarna een sprint). Morgen: langzaam
eten, en de maaltijd laten zakken voor het rennen.
\end{solution}

\begin{solution}{exo:g4:journey-of-food:10}
Doorgeslikt voedsel zit nog in de buis --- als het gat van een donut,
nog niet echt in het lichaam. Het moet eerst tot \omterm{def:g4:journey-of-food:digestion}{voedingsstoffen}
verteerd worden; pas in de \omterm{def:g4:journey-of-food:tract}{dunne darm} steken die de wand over naar het
\omterm{prop:g4:journey-of-food:absorb}{bloed}. De \omterm{def:g4:journey-of-food:tract}{mond} binnengaan is niet het lichaam binnengaan.
\end{solution}

\begin{solution}{exo:g4:journey-of-food:11}
De \omterm{def:g4:journey-of-food:digestion}{vertering} moet voedsel uit elkaar halen tot \omterm{def:g4:journey-of-food:digestion}{voedingsstoffen}, en gras
is veel moeilijker uit elkaar te halen dan brood --- taaie vezels,
weinig voedingswaarde per hap
(\cref{prop:g2:what-animals-eat:daily}). Dus brengt de graseter
zwaardere machinerie mee: maalkiezen die twee keer gebruikt worden (het
herkauwen), vier maagkamers voor een lange bewerking, en darmhelpers
die kunnen afbreken wat haar eigen sappen niet kunnen. Taaiere invoer,
grotere fabriek.
\end{solution}
